% ============================================================
%  §4.1.4  LLM Training
%  Drop-in replacement for the existing \subsubsection block.
%  Requires: \usepackage{natbib}   (add to preamble)
%            \usepackage{booktabs} (for \toprule etc. in tables)
% ============================================================

\subsubsection{LLM Training}

\paragraph{QLoRA.}

The Llama~3.1 8B model is built on the transformer architecture, which
processes text through a deep stack of identical layers, each comprising a
multi-head self-attention mechanism and a position-wise feed-forward network
(MLP). Together, these components learn rich contextual representations of
language through weight matrices distributed across 32 such layers and totalling
eight billion parameters. Full fine-tuning of a model at this scale requires
updating every parameter simultaneously, which demands upwards of 60~GB of GPU
memory in 16-bit precision---well beyond the capacity of consumer hardware.

Low-Rank Adaptation \citep[LoRA;][]{hu2021lora} addresses this constraint by
freezing all pre-trained weights and injecting pairs of small trainable matrices
into targeted layers of the transformer. For a frozen weight matrix
$W_0 \in \mathbb{R}^{d \times k}$, LoRA parameterises the incremental update as
$\Delta W = BA$, where $B \in \mathbb{R}^{d \times r}$ and
$A \in \mathbb{R}^{r \times k}$ with rank $r \ll \min(d, k)$.
Only $A$ and $B$ are trained; the adapted forward pass becomes
$h = W_0 x + \tfrac{\alpha}{r} BA x$,
where $\alpha$ is a scaling hyperparameter controlling the relative magnitude of
the low-rank update. This reduces the number of trainable parameters to under
one per cent of the full model while achieving performance comparable to full
fine-tuning on classification tasks \citep{hu2021lora}.

Quantized LoRA \citep[QLoRA;][]{dettmers2023qlora} extends this approach by
additionally compressing the frozen base-model weights to 4-bit NormalFloat
(NF4), a data type that is information-theoretically optimal for
normally-distributed weights. This reduces the memory footprint of the base
model by approximately $4\times$ relative to 16-bit precision, making it
feasible to fine-tune an 8B-parameter model on a single consumer GPU.
Crucially, quantization is applied only to the frozen weights; the LoRA
adapters $A$ and $B$ are maintained in full 16-bit precision throughout
training, preserving gradient fidelity. QLoRA further employs double
quantization---compressing the quantization constants themselves---and paged
optimisers to manage memory spikes during gradient checkpointing.
\citet{dettmers2023qlora} demonstrate that the reduction in representational
flexibility introduced by 4-bit quantization is best compensated by injecting
LoRA adapters across \emph{all} linear layers of the transformer, rather than
restricting injection to the attention matrices alone. This all-layer strategy
is adopted here, and carries a further consequence: final task performance
becomes largely insensitive to the choice of rank $r$, shifting the primary
regularisation burden to the remaining hyperparameters
\citep{dettmers2023qlora}.

\paragraph{Hyperparameter Optimisation.}

Five hyperparameters govern the expressivity and regularisation of QLoRA
fine-tuning: rank $r$, scaling multiplier $m$ (where $\alpha = m \cdot r$),
LoRA dropout, learning rate, and AdamW weight decay. Because the training
corpus comprises approximately 700 annotated sentences in the initial active
learning round, overfitting is the dominant risk; optimal values are not
directly transferable from the large-dataset regimes studied in the bulk of the
fine-tuning literature and must be determined empirically. Preliminary review of
comparable financial NLP fine-tuning experiments on Llama~3.1 8B was used to
establish principled bounds for each parameter, and systematic search was then
performed using Optuna \citep{akiba2019optuna} with Tree-structured Parzen
Estimator (TPE) sampling over 50 trials.

\textit{Rank $r$.}
The rank controls the dimensionality of the LoRA update matrices and, by
extension, the number of trainable parameters introduced per layer. In a
low-data regime the primary concern is that higher rank increases overfitting
risk, a consideration that outweighs the marginal benefit of additional
capacity. The closest analogue to this study in the literature is FinLoRA
\citep{finlora2024}, which applies 4-bit QLoRA fine-tuning of Llama~3.1 8B to
financial classification tasks including sentiment analysis, and reports that
rank 4 achieves performance comparable to rank 8 while requiring less memory.
The search space $r \in \{2, 4, 8\}$ therefore spans from a minimally
expressive adapter to the FinLoRA upper bound, with $r = 4$ as the informed
prior.

\textit{Alpha multiplier $m$.}
The scaling parameter $\alpha$ is parameterised as $m \cdot r$ with
$m \in \{1, 2\}$. A multiplier of 2 is empirically optimal across most tasks
\citep{raschka2023lora}, and tuning $m$ jointly with the learning rate avoids
confounding the two, since \citet{hu2021lora} note that adjusting $\alpha$ is
broadly analogous to rescaling the learning rate.

\textit{LoRA dropout.}
Dropout applied to the adapter layers provides the primary regularisation
channel. \citet{dettmers2023qlora} recommend a rate of 0.05 for 7--8B models,
but this baseline was established on large instruction datasets.
In the low-data regime of this study, stronger regularisation is warranted, and
the search space $p \in \{0.05, 0.1, 0.2\}$ extends upward from the Dettmers
baseline accordingly.

\textit{Learning rate.}
The learning rate is the most sensitive QLoRA hyperparameter, with small
differences producing large swings in task performance \citep{biderman2024lora}.
It is searched on a log scale over $[10^{-4}, 5 \times 10^{-4}]$, with
$2 \times 10^{-4}$ as the standard starting point in the literature
\citep{dettmers2023qlora, finlora2024}; the upper bound is included to allow the
search to confirm it is suboptimal under this data regime.

\textit{Weight decay.}
\citet{hu2021lora} note that applying weight decay to the LoRA matrices $A$ and
$B$ implicitly decays updates back toward the pre-trained weights, providing a
natural guard against catastrophic forgetting that is distinct from dropout.
The search space $\lambda \in \{0.01, 0.05, 0.1\}$ treats weight decay as a
second, independent regularisation axis alongside dropout.

To avoid wasteful evaluation of poor configurations on a constrained GPU budget,
the Optuna search employs a Median Pruner: the first 10 trials perform unconstrained
random exploration to establish a performance baseline, after which any trial
whose intermediate validation score falls below the median of completed trials at
the same step (evaluated after a 50-step warm-up) is terminated early.
The full search space is summarised in Table~\ref{tab:optuna_search}.

\begin{table}[ht]
\centering
\caption{Optuna hyperparameter search space for QLoRA fine-tuning of Llama~3.1 8B.
         50 trials with TPE sampling and Median Pruner (10 startup trials,
         50-step warm-up).}
\label{tab:optuna_search}
\begin{tabular}{p{3.2cm} p{3.6cm} p{4.8cm} p{2.5cm}}
\toprule
\textbf{Parameter} & \textbf{Search space} & \textbf{Rationale} & \textbf{Source} \\
\midrule
Rank $r$
  & $\{2,\, 4,\, 8\}$
  & $r=4$ optimal for 4-bit QLoRA on financial classification; $r=8$ upper bound
  & \citet{finlora2024} \\[4pt]
Alpha multiplier $m$
  & $\{1,\, 2\}$
  & $m=2$ empirically optimal; interacts with LR
  & \citet{raschka2023lora} \\[4pt]
LoRA dropout
  & $\{0.05,\, 0.1,\, 0.2\}$
  & 0.05 Dettmers baseline; extended upward for small-data regime
  & \citet{dettmers2023qlora} \\[4pt]
Learning rate
  & $[10^{-4},\; 5{\times}10^{-4}]$ (log)
  & $2{\times}10^{-4}$ standard; upper bound confirms small-data preference
  & \citet{dettmers2023qlora, finlora2024} \\[4pt]
Weight decay $\lambda$
  & $\{0.01,\, 0.05,\, 0.1\}$
  & Second regularisation axis; decays updates toward pre-trained weights
  & \citet{hu2021lora} \\
\bottomrule
\end{tabular}
\end{table}

\paragraph{Training Configuration.}

Beyond the Optuna-optimised hyperparameters, a set of fixed training parameters
governs the mechanics of the fine-tuning procedure. These are chosen to ensure
stable convergence within the memory constraints of a single consumer GPU
(NVIDIA RTX~4070 Ti Super, 16~GB VRAM) and to guard against overfitting.

Training proceeds for a maximum of 3 epochs. Although broader fine-tuning
practice sometimes employs up to 5 epochs, the small training set of
approximately 700 sentences places this study in a regime where the model risks
memorising the training examples rather than generalising; the conservative
ceiling of 3 epochs is motivated by this constraint
\citep{dettmers2023qlora, finlora2024}. The checkpoint maximising macro-averaged
F1 on the validation set is retained at the end of training via early stopping,
providing a further safeguard against epoch-level overfitting. Macro-F1 is
chosen as the selection metric for robustness to class imbalance across the seven
output fields of the annotation schema.

The effective batch size is set to 16, achieved via a per-device batch size of 8
with 2 gradient accumulation steps, which sits within the range supported by
16~GB VRAM under 4-bit quantization. A cosine learning rate schedule with a
warmup ratio of 0.1 decays the learning rate smoothly from its initial value to
near zero, reducing the risk of large gradient updates disrupting pre-trained
representations in the final steps \citep{raschka2023lora}. Training is
performed in BFloat16 (bf16) precision, which is natively supported by the
RTX~4070 Ti Super architecture and provides the numerical range of 32-bit
floating point at 16-bit memory cost, making it preferable to FP16 for training
stability. The paged AdamW 8-bit optimiser, standard in QLoRA pipelines
\citep{dettmers2023qlora}, is used to prevent memory spikes during gradient
checkpointing. The full training configuration is reported in
Table~\ref{tab:training_config}.

\begin{table}[ht]
\centering
\caption{Fixed training configuration for QLoRA fine-tuning of Llama~3.1 8B.}
\label{tab:training_config}
\begin{tabular}{p{4.5cm} p{3.5cm} p{6.5cm}}
\toprule
\textbf{Parameter} & \textbf{Value} & \textbf{Rationale} \\
\midrule
Max epochs              & 3
  & Conservative ceiling for $n \approx 700$; early stopping on val.\ F1 \\
Per-device batch size   & 8
  & Maximum safe for 16~GB VRAM with 4-bit QLoRA \\
Gradient accumulation   & 2 steps (eff.\ batch $= 16$)
  & Simulates larger batch without additional memory \\
Eval.\ \& save strategy & Per epoch
  & Enables checkpoint selection on validation F1 \\
Best model metric       & Macro F1
  & Robust to class imbalance across output fields \\
LR scheduler            & Cosine
  & Smooth decay to near zero; standard for LoRA \\
Warmup ratio            & 0.1
  & Stabilises early training; prevents large initial updates \\
Precision               & bf16
  & Native RTX~4070 Ti Super support; more stable than fp16 \\
Optimiser               & Paged AdamW 8-bit
  & QLoRA standard; manages memory spikes during checkpointing \\
\bottomrule
\end{tabular}
\end{table}
